\section{Conclusion}

Adaptation burden under drift is not organized by a single memory coordinate.
The supported empirical structure is hierarchical. Across families, effective
spectral rank identifies degradation regimes. Within a recurrent local regime,
activation velocity controls burden more strongly than rank even when rank still
varies broadly. That separation explains why static endpoint metrics are blind to
temporal adaptation quality: regime selection and within-regime scale control
move burden through different channels.

The next theorem-level empirical step is not to enlarge the current claim, but to
stabilize the open mechanism. Architecture-moderated rank effects and broader
within-family burden laws remain open problems built on top of the hierarchical
structure established here.
